\documentclass[10pt,a4paper]{amsart}
\usepackage[margin=19mm]{geometry}
\usepackage{amsmath,amssymb,mathtools}
\usepackage[scr=rsfs,cal=euler]{mathalfa}
\usepackage{xfrac}
\usepackage[unicode,hidelinks,bookmarks=false]{hyperref}
\newcommand{\ie}{i.e.\ }
\newcommand{\cf}{cf.\ }
\newcommand{\resp}{respectively\ }
\newcommand{\Subsection}{\subsection}
\providecommand{\nobreakdash}{\nobreak\hbox{-}}
\setlength{\emergencystretch}{2em}
\multlinegap0pt
\usepackage{bm}
\usepackage{bbm}
\usepackage{dsfont}
\usepackage{xspace}
\usepackage{cancel}
\usepackage{mathtools}
\usepackage{amsthm}
\makeatletter
\@ifundefined{theorem}{\newtheorem{theorem}{Theorem}}{}
\@ifundefined{lemma}{\newtheorem{lemma}[theorem]{Lemma}}{}
\makeatother
\title[The Frog Model on $\mathbb Z$ with Random Discrete Weibull L]{The Frog Model on $\mathbb Z$ with Random Discrete Weibull Lifetimes and Biased Nearest-Neighbour Random Walks}
\author{Ram\'irez-Gonz\'alez J.H., Prates Machado Fabio}
\date{Source: arXiv:2606.11068v1 (CC BY 4.0)}
\begin{document}
\maketitle

\noindent\emph{Source and licence.} The Frog Model on $\mathbb Z$ with Random Discrete Weibull Lifetimes and Biased Nearest-Neighbour Random Walks. Version arXiv:2606.11068v1 (CC BY 4.0), distributed under the Creative Commons Attribution 4.0 International licence, as recorded in the arXiv licence metadata for this exact version and shown on the abstract page https://arxiv.org/abs/2606.11068v1. Full rights notice: licenses/arxiv-2606.11068-CC-BY-4.0.txt.\par\smallskip

\noindent\textit{Editorial scope.} Counted unit: the Lemma whose source label is \texttt{lem:drifted-hitting-time} in the article (source0.tex lines 731--750, source label \texttt{lem:drifted-hitting-time}), delivered together with the article's own complete original proof of it (source0.tex lines 752--983, one \texttt{proof} environment of 232 lines). No mathematics is written, repaired or re-derived: statement and proof are the article's own text line for line. Required context retained uncounted: none. Every definition, display and label of those lines is retained unchanged. Omitted from this package and not redistributed: 1--730, 751--751, 984--1594. Nothing inside a retained excerpt is omitted or altered: every retained body is delivered exactly as the article's lines are written, so no omission receipt of any kind accompanies this package. Citations: the retained excerpts carry no citation command at all, so no bibliography entry of the article is transcribed and no reference is redistributed. Reference adaptation: none was needed and none was made; every reference inside the delivered unit resolves inside it, so no unresolved reference or citation is produced. Printed slips kept literal: no printed word, symbol, hypothesis or number of the counted statement or of its proof was altered. Layout and class: the article's own class is \texttt{amsart} with packages including inputenc, fontenc, amsmath, amssymb, amsfonts, amsthm, geometry, enumitem, xcolor, hyperref, natbib, url, lineno, graphicx; the wrapper is the lane's pinned \texttt{amsart} editorial wrapper at 10pt on A4, which restates the theorem-like environments the retained text needs and adds the article's own macro definitions that the retained text needs (none), with extra wrapper packages (bm, bbm, dsfont, xspace, cancel, mathtools). The delivered numbering is the unit's own. Assets: no retained span contains a figure environment, a graphics-inclusion command, a PDF-page inclusion, a table or a TikZ picture; no image file is copied into this package, the package declares no asset dependency and no image file is redistributed with it. Licence: version arXiv:2606.11068v1 of this article is distributed under the Creative Commons Attribution 4.0 International licence, as recorded in the arXiv licence metadata for this exact version and shown on the abstract page https://arxiv.org/abs/2606.11068v1. A full rights notice is delivered at licenses/arxiv-2606.11068-CC-BY-4.0.txt. Characters: every retained excerpt is pure ASCII, so the delivered engine, XeLaTeX with its default Latin Modern text font, reports no missing character.
\begin{lemma}\label{lem:drifted-hitting-time}
The following hitting-time asymptotics hold:
\[
  \frac{\tau_n^*}{n}
  \xrightarrow[n\to\infty]{a.s.}
  \frac1v.
\]
Moreover, if \(r>1/2\), then
\[
  \frac{\tau_n^+}{n}
  \xrightarrow[n\to\infty]{a.s.}
  \frac1{2r-1},
\]
whereas, if \(r<1/2\), then
\[
  \frac{\tau_n^-}{n}
  \xrightarrow[n\to\infty]{a.s.}
  \frac1{1-2r}.
\]
\end{lemma}

\begin{proof}
It is enough to give the argument for \(r>1/2\), and then use reflection for
the case \(r<1/2\). Assume first that \(r>1/2\), and set \(  \alpha=2r-1>0\). By the strong law of large numbers,
\begin{equation}\label{eq:slln-drift-positive}
  \frac{S_m}{m}\longrightarrow \alpha
  \qquad\text{a.s.}
\end{equation}
We work on an event of probability one on which \eqref{eq:slln-drift-positive}
holds.

We first prove the asymptotic for \(\tau_n^+\). Let \(b>1/\alpha\). Choose
\(\varepsilon>0\) such that \(b(\alpha-\varepsilon)>1\). By
\eqref{eq:slln-drift-positive}, there exists \(m_0=m_0(\varepsilon)\) such that
\[
  S_m\ge (\alpha-\varepsilon)m,
  \qquad m\ge m_0 .
\]
Put \(m_n:=\lfloor bn\rfloor\). For all sufficiently large \(n\),
\(m_n\ge m_0\), and
\[
  \frac{(\alpha-\varepsilon)m_n}{n}
  \longrightarrow
  b(\alpha-\varepsilon)>1.
\]
Hence \(S_{m_n}\ge(\alpha-\varepsilon)m_n>n\) for all sufficiently large \(n\).
Since the walk is nearest-neighbour and starts from \(0\), it must have visited
\(n\) by time \(m_n\). Therefore \(\tau_n^+\le m_n\) eventually, and
\[
  \limsup_{n\to\infty}\frac{\tau_n^+}{n}\le b.
\]
Letting \(b\downarrow1/\alpha\), we obtain
\begin{equation}\label{eq:tau-plus-limsup}
  \limsup_{n\to\infty}\frac{\tau_n^+}{n}
  \le
  \frac1\alpha .
\end{equation}

For the lower bound, let \(a<1/\alpha\). If \(a\le1\), then
\(\tau_n^+\ge n>an\) for every \(n\). Assume therefore that \(a>1\), and choose
\(\varepsilon>0\) such that
\[
  \alpha+\varepsilon<\frac1a.
\]
Again by \eqref{eq:slln-drift-positive}, there exists
\(m_1=m_1(\varepsilon)\) such that
\[
  \frac{S_m}{m}\le \alpha+\varepsilon,
  \qquad m\ge m_1 .
\]
Suppose that, along an infinite subsequence \((n_j^{(1)})_{j\ge1}\),
\[
  \tau_{n_j^{(1)}}^+\le a n_j^{(1)} .
\]
Since the walk is nearest-neighbour, \(\tau_{n_j^{(1)}}^+\ge n_j^{(1)}\), and
therefore \(\tau_{n_j^{(1)}}^+\to\infty\). Thus, after discarding finitely many
indices, \(\tau_{n_j^{(1)}}^+\ge m_1\). For such \(j\),
\[
  S_{\tau_{n_j^{(1)}}^+}=n_j^{(1)},
\]
and hence
\[
  \frac{S_{\tau_{n_j^{(1)}}^+}}{\tau_{n_j^{(1)}}^+}
  =
  \frac{n_j^{(1)}}{\tau_{n_j^{(1)}}^+}
  \ge
  \frac1a
  >
  \alpha+\varepsilon,
\]
which contradicts the definition of \(m_1\). Hence \(\tau_n^+>an\) eventually,
and
\[
  \liminf_{n\to\infty}\frac{\tau_n^+}{n}\ge a.
\]
Letting \(a\uparrow1/\alpha\), we get
\begin{equation}\label{eq:tau-plus-liminf}
  \liminf_{n\to\infty}\frac{\tau_n^+}{n}
  \ge
  \frac1\alpha .
\end{equation}
Combining \eqref{eq:tau-plus-limsup} and \eqref{eq:tau-plus-liminf} gives
\[
  \frac{\tau_n^+}{n}
  \longrightarrow
  \frac1\alpha
  =
  \frac1{2r-1}.
\]

We now prove the corresponding assertion for \(\tau_n^*\). Since
\(\tau_n^*\le\tau_n^+\), the upper bound gives
\[
  \limsup_{n\to\infty}\frac{\tau_n^*}{n}
  \le
  \frac1\alpha.
\]
For the lower bound, let \(a<1/\alpha\). If \(a\le1\), then
\(\tau_n^*\ge n>an\). Assume \(a>1\), and choose \(\varepsilon>0\) such that
\[
  \alpha+\varepsilon<\frac1a,
  \qquad
  \alpha-\varepsilon>0.
\]
By \eqref{eq:slln-drift-positive}, there exists \(m_2=m_2(\varepsilon)\) such
that
\[
  \alpha-\varepsilon
  \le
  \frac{S_m}{m}
  \le
  \alpha+\varepsilon,
  \qquad m\ge m_2 .
\]
Suppose that, along an infinite subsequence \((n_j^{(2)})_{j\ge1}\),
\[
  \tau_{n_j^{(2)}}^*\le a n_j^{(2)} .
\]
Since \(\tau_{n_j^{(2)}}^*\ge n_j^{(2)}\), we have
\(\tau_{n_j^{(2)}}^*\to\infty\). After discarding finitely many terms, assume
\[
  \tau_{n_j^{(2)}}^*\ge m_2,
  \qquad j\ge1.
\]
For each \(j\),
\[
  S_{\tau_{n_j^{(2)}}^*}=n_j^{(2)}
  \qquad\text{or}\qquad
  S_{\tau_{n_j^{(2)}}^*}=-n_j^{(2)}.
\]
Define
\[
  J_+:=\left\{
    j\ge1:
    S_{\tau_{n_j^{(2)}}^*}=n_j^{(2)}
  \right\},
  \qquad
  J_-:=\left\{
    j\ge1:
    S_{\tau_{n_j^{(2)}}^*}=-n_j^{(2)}
  \right\}.
\]
Since \(\mathbb N=J_+\cup J_-\), at least one of \(J_+\) and \(J_-\) is
infinite.

If \(J_+\) is infinite, take a further infinite subsequence
\((n_\ell^{(2,+)})_{\ell\ge1}\). Then
\[
  \frac{S_{\tau_{n_\ell^{(2,+)}}^*}}{\tau_{n_\ell^{(2,+)}}^*}
  =
  \frac{n_\ell^{(2,+)}}{\tau_{n_\ell^{(2,+)}}^*}
  \ge
  \frac1a
  >
  \alpha+\varepsilon,
\]
contradicting the upper bound \(S_m/m\le\alpha+\varepsilon\) for \(m\ge m_2\).
If \(J_-\) is infinite, take a further infinite subsequence
\((n_\ell^{(2,-)})_{\ell\ge1}\). Then
\[
  \frac{S_{\tau_{n_\ell^{(2,-)}}^*}}{\tau_{n_\ell^{(2,-)}}^*}
  =
  -\frac{n_\ell^{(2,-)}}{\tau_{n_\ell^{(2,-)}}^*}
  <0,
\]
contradicting the lower bound \(S_m/m\ge\alpha-\varepsilon>0\) for
\(m\ge m_2\). Both cases are impossible. Therefore \(\tau_n^*>an\)
eventually, and
\[
  \liminf_{n\to\infty}\frac{\tau_n^*}{n}\ge a.
\]
Letting \(a\uparrow1/\alpha\), we obtain
\[
  \liminf_{n\to\infty}\frac{\tau_n^*}{n}
  \ge
  \frac1\alpha.
\]
Together with the upper bound, this proves
\[
  \frac{\tau_n^*}{n}
  \longrightarrow
  \frac1\alpha
  =
  \frac1{2r-1}
  =
  \frac1v .
\]

It remains to treat \(r<1/2\). Define \(\widetilde S_m:=-S_m\). Then
\((\widetilde S_m)_{m\ge0}\) is a nearest-neighbour random walk with
\[
  \mathbb P(\widetilde X_1=1)=1-r>\frac12,
  \qquad
  \mathbb P(\widetilde X_1=-1)=r,
\]
and drift \(1-2r=v>0\). Applying the previous part to
\((\widetilde S_m)_{m\ge0}\) yields
\[
  \frac1n\inf\{m\ge1:\widetilde S_m=n\}
  \longrightarrow
  \frac1{1-2r}.
\]
Since
\[
  \inf\{m\ge1:\widetilde S_m=n\}
  =
  \inf\{m\ge1:S_m=-n\}
  =
  \tau_n^-,
\]
we get
\[
  \frac{\tau_n^-}{n}
  \longrightarrow
  \frac1{1-2r}.
\]
Moreover,
\[
  \inf\{m\ge1:|\widetilde S_m|=n\}
  =
  \inf\{m\ge1:|S_m|=n\}
  =
  \tau_n^*,
\]
and therefore
\[
  \frac{\tau_n^*}{n}
  \longrightarrow
  \frac1{1-2r}
  =
  \frac1v .
\]
\end{proof}

\end{document}
